\documentclass[12pt,a4paper]{article}

% =====================================================
% Paquetes
% =====================================================
\usepackage[utf8]{inputenc}
\usepackage[T1]{fontenc}
\usepackage[spanish,es-nodecimaldot]{babel}
\usepackage{lmodern}
\usepackage{geometry}
\usepackage{setspace}
\usepackage{graphicx}
\usepackage{booktabs}
\usepackage{array}
\usepackage{xcolor}
\usepackage{listings}
\usepackage{microtype}
\usepackage{hyperref}
\usepackage{enumitem}
\usepackage{booktabs}
\usepackage{array}
\usepackage{float}
\usepackage{longtable}
\usepackage{graphicx}
\usepackage{subcaption}


% =====================================================
% Configuración de página y estilo
% =====================================================
\geometry{top=2.5cm,bottom=2.5cm,left=3cm,right=3cm}
\onehalfspacing
\setlength{\parindent}{1.2em}
\setlength{\parskip}{0.35em}

\definecolor{azul}{RGB}{31,78,121}
\definecolor{grisclaro}{RGB}{245,247,250}

\hypersetup{
    colorlinks=true,
    linkcolor=azul,
    citecolor=azul,
    urlcolor=azul,
    pdftitle={Recocido Simulado para la resolución del problema del agente viajero},
    pdfauthor={Ramírez Juárez María Fernanda}
}

\lstset{
    basicstyle=\ttfamily\small,
    backgroundcolor=\color{grisclaro},
    frame=single,
    rulecolor=\color{black!15},
    columns=fullflexible,
    keepspaces=true,
    showstringspaces=false,
    breaklines=true,
    tabsize=2
}

\setlist[description]{style=nextline,leftmargin=0pt,labelindent=0pt,itemsep=0.6em}

% =====================================================
% Información del documento
% =====================================================
\title{\textbf{La heurística Old Bachelor Acceptance (OBA) para la resolución de Job Shop Scheduling Problem }}
\author{Ramírez Juárez María Fernanda}
\date{\today}

% =====================================================
% Documento
% =====================================================
\begin{document}

\maketitle

% =====================================================
% 1. INTRODUCCIÓN
% =====================================================
\section{Introducción}
\label{sec:introduccion}

El problema de programación de trabajos en máquinas, conocido como
\textit{Job-Shop Scheduling Problem} (JSSP), es un problema clásico
de optimización combinatoria que consiste en determinar el orden
en que un conjunto de trabajos debe ser procesado por diferentes
máquinas. Cada trabajo está compuesto por una secuencia
predeterminada de operaciones, donde cada operación requiere una
máquina específica y un tiempo de procesamiento determinado.
El objetivo principal es encontrar una programación que permita
completar todos los trabajos en el menor tiempo posible, respetando
las restricciones de precedencia entre operaciones y la disponibilidad
de los recursos.

El JSSP presenta una dificultad considerable debido a la gran
cantidad de posibles combinaciones en las que pueden organizarse
las operaciones. Conforme aumenta el número de trabajos, máquinas
y operaciones, el espacio de búsqueda crece significativamente,
haciendo que encontrar soluciones óptimas resulte computacionalmente
costoso. Por esta razón, las heurísticas de optimización
combinatoria representan una alternativa para obtener soluciones
de buena calidad en tiempos de ejecución razonables.

En este proyecto se propone utilizar la heurística
\textit{Old Bachelor Acceptance} (OBA), específicamente su variante
OBA2-N, para abordar el problema JSSP. OBA es un método de búsqueda
local basado en aceptación por umbrales (\textit{Threshold
Accepting}), propuesto por Huy, Kahng y Tsao. Su principal
característica es el uso de un umbral dinámico y no monótono que
permite controlar la aceptación de soluciones que pueden representar
un empeoramiento respecto a la solución actual. De esta manera,
el algoritmo busca evitar quedar atrapado en óptimos locales y
mantener la exploración del espacio de soluciones.

La variante OBA2 utiliza una estrategia que combina movimientos
de descenso pronunciado con incrementos moderados del umbral.
Esto permite intensificar la búsqueda cuando se encuentran
soluciones favorables y relajar gradualmente el criterio de
aceptación cuando se presentan dificultades para encontrar nuevos
movimientos aceptables. Por su parte, OBA2-N introduce una
modificación en el criterio de aceptación: los valores negativos
del umbral se consideran como cero al evaluar si una nueva
solución puede ser aceptada. Esta modificación permite estudiar
el comportamiento de la búsqueda sin utilizar umbrales negativos
para aceptar movimientos.

Finalmente, el propósito de este proyecto es implementar y evaluar
la heurística OBA2-N como método de búsqueda para resolver
instancias del problema JSSP, utilizando como criterio de
optimización la minimización del tiempo total de finalización
de los trabajos, conocido como \textit{makespan}. A partir de
los resultados obtenidos, se busca analizar el comportamiento
de la heurística y su capacidad para encontrar programaciones
de buena calidad.

% =====================================================
% 2. DESARROLLO
% =====================================================
\section{Desarrollo}

Para resolver el problema de calendarización de trabajos en máquinas
(Job-Shop Scheduling Problem, JSSP), se desarrolló una aplicación en
el lenguaje de programación Go que utiliza la heurística
\textit{Old Bachelor Acceptance}, específicamente su variante
OBA2-N. El desarrollo se organizó mediante módulos independientes,
con el propósito de separar la definición del problema, la
representación de las soluciones, su evaluación y el proceso de
búsqueda.

El procedimiento general consiste en leer una instancia del JSSP,
construir una solución inicial, evaluar su calidad y explorar
soluciones vecinas mediante movimientos sobre el orden de
procesamiento de las operaciones. La heurística OBA2-N determina
si cada solución candidata es aceptada, utilizando un umbral
dinámico. Finalmente, se registran las mejores soluciones y las
métricas obtenidas durante las ejecuciones.

\subsection{Modelado del problema}

El primer paso consistió en representar formalmente las instancias
del JSSP. Una instancia está compuesta por un conjunto de trabajos,
donde cada trabajo contiene una secuencia ordenada de operaciones.
Cada operación requiere una máquina específica y posee un tiempo
de procesamiento determinado.

Para representar esta información se definieron tres estructuras
principales:

\begin{itemize}
    \item \textbf{Instance:} almacena las dimensiones del problema,
    el conjunto de trabajos, las operaciones y las cotas inferior
    y superior.

    \item \textbf{Job:} representa un trabajo y contiene su
    identificador y la secuencia ordenada de operaciones que
    deben ejecutarse.

    \item \textbf{Operation:} representa una operación mediante
    su identificador, el trabajo al que pertenece, su posición
    dentro del trabajo, la máquina requerida y su duración.
\end{itemize}

Esta representación permite identificar las operaciones y
consultar sus características durante la construcción y evaluación
de los calendarios.

\subsubsection{Lectura y validación de instancias}

Se desarrolló un mecanismo para leer las instancias desde archivos
de texto. El formato de entrada utilizado comienza con dos valores
enteros que indican el número de trabajos y el número de máquinas.
Posteriormente, se especifican los pares máquina-duración
correspondientes a cada operación.

De esta manera, cada par representa la máquina que debe procesar
una operación y el tiempo necesario para completarla. El parser
interpreta estos datos y construye las estructuras de trabajos y
operaciones utilizadas por el resto del programa.

Después de la lectura se realiza una validación de la instancia,
comprobando aspectos como:

\begin{itemize}
    \item Que las dimensiones de la instancia sean válidas.
    \item Que las operaciones tengan identificadores consistentes.
    \item Que las máquinas especificadas pertenezcan al rango
    permitido.
    \item Que las duraciones de procesamiento sean positivas.
    \item Que la cantidad de operaciones corresponda con las
    dimensiones de la instancia.
\end{itemize}

Esta etapa permite detectar entradas incorrectas antes de iniciar
el proceso de búsqueda.

\subsubsection{Cálculo de cotas}

Para evaluar las soluciones se establecieron una cota inferior
(\textit{Lower Bound}, LB) y una cota superior
(\textit{Upper Bound}, UB).

La cota superior se obtiene sumando las duraciones de todas las
operaciones de la instancia:

\[
UB = \sum_{i=1}^{n}\sum_{j=1}^{m} t(O_{ij})
\]

donde $t(O_{ij})$ representa la duración de la operación $j$
del trabajo $i$.

Por otra parte, la cota inferior se calcula considerando dos
restricciones del problema: la carga de trabajo de cada trabajo
y la carga de procesamiento de cada máquina. Por lo tanto:

\[
LB = \max(LB_{\text{trabajos}},
          LB_{\text{máquinas}})
\]

La primera corresponde a la mayor suma de duraciones entre
todos los trabajos, mientras que la segunda representa la mayor
suma de duraciones entre todas las máquinas.

Estas cotas se utilizan posteriormente para normalizar el valor
de la función objetivo.

\subsection{Representación de las soluciones}

Una vez definido el modelo del problema, se estableció una
representación para los calendarios de procesamiento. Se utilizó
una estructura denominada \texttt{Schedule}, basada en una lista
de operaciones para cada máquina.

La estructura \texttt{MachineOrder} almacena los identificadores
de las operaciones en el orden en que serán procesadas por cada
máquina. Por consiguiente, cada posición de la estructura
corresponde a una máquina y contiene la secuencia de operaciones
asignadas a ella.

Esta representación permite modificar el orden de procesamiento
de las operaciones sin alterar la máquina a la que pertenece
cada una.

Es importante señalar que agrupar las operaciones por máquina
no garantiza por sí mismo que el calendario sea factible. También
deben respetarse las restricciones de precedencia entre las
operaciones de un mismo trabajo. Por esta razón, la factibilidad
se verifica durante la evaluación de cada solución.

\subsection{Evaluación de soluciones}

La evaluación de una solución es una etapa fundamental, ya que
permite determinar su calidad y comprobar si cumple las
restricciones del JSSP.

Para realizar esta evaluación se construye un grafo dirigido
que representa las relaciones de precedencia entre las
operaciones. Se consideran dos tipos de restricciones:

\begin{itemize}
    \item \textbf{Precedencia entre operaciones de un trabajo:}
    las operaciones deben ejecutarse respetando el orden
    establecido en cada trabajo.

    \item \textbf{Orden de procesamiento de una máquina:}
    las operaciones asignadas a una misma máquina deben
    procesarse siguiendo el orden definido en el calendario.
\end{itemize}

Cada operación se representa mediante un vértice del grafo,
mientras que las restricciones de precedencia se representan
mediante aristas dirigidas.

\subsubsection{Detección de ciclos y factibilidad}

Para comprobar si un calendario es factible se utiliza el
ordenamiento topológico mediante el algoritmo de Kahn.

Este procedimiento permite procesar los vértices del grafo
respetando sus dependencias. Si se consigue procesar la totalidad
de las operaciones, significa que el grafo es acíclico y que
el calendario satisface las restricciones de precedencia
modeladas.

En cambio, si quedan vértices sin procesar, se detecta la
existencia de ciclos o dependencias que impiden completar
el ordenamiento topológico. En este caso, la solución se
considera infactible.

\subsubsection{Cálculo del makespan}

Cuando el grafo es acíclico, se calculan los tiempos de inicio
y finalización de las operaciones respetando las dependencias
establecidas.

El objetivo principal del JSSP es minimizar el tiempo total
necesario para completar todos los trabajos, conocido como
\textit{makespan}. Este valor corresponde al mayor tiempo
de finalización entre todas las operaciones:

\[
C_{\max} = \max_{i,j} C_{ij}
\]

donde $C_{ij}$ representa el tiempo de finalización de la
operación $j$ del trabajo $i$.

Por lo tanto, una solución con un menor valor de
$C_{\max}$ representa un calendario de procesamiento más
eficiente, siempre que sea factible.

\subsubsection{Función objetivo}

Para evaluar tanto la calidad como la factibilidad de las
soluciones se utiliza una función de costo que combina el
makespan normalizado con una penalización por ciclos:

\[
f(X) =
\frac{C_{\max}(X)-LB}{UB-LB}
+
LB \cdot V(X)
\]

donde:

\begin{itemize}
    \item $X$ representa el calendario evaluado.
    \item $C_{\max}(X)$ es el makespan de la solución.
    \item $LB$ y $UB$ representan las cotas inferior y superior.
    \item $V(X)$ representa la cantidad de vértices que quedaron
    sin procesar durante la evaluación topológica.
\end{itemize}

El primer término permite normalizar el makespan con respecto
a las cotas calculadas para la instancia. El segundo término
penaliza las soluciones que presentan ciclos.

De esta manera, las soluciones factibles se evalúan mediante
su makespan normalizado, mientras que las soluciones
infactibles reciben una penalización que incrementa su costo.

\subsection{Generación de soluciones vecinas}

Para explorar diferentes calendarios se implementó un mecanismo
de generación de soluciones vecinas. La vecindad se basa en
intercambiar operaciones adyacentes dentro de la secuencia
correspondiente a una máquina.

Si una máquina contiene una secuencia de operaciones como:

\[
M_k = [O_1,O_2,O_3,O_4]
\]

un movimiento puede intercambiar dos operaciones consecutivas,
obteniendo, por ejemplo:

\[
M_k' = [O_1,O_3,O_2,O_4]
\]

Este movimiento modifica el orden de procesamiento de las
operaciones sin cambiar su asignación a las máquinas.

Cada movimiento se aplica sobre una copia del calendario
original, de manera que la solución actual no sea modificada
accidentalmente durante la generación de candidatos.

Después de realizar el intercambio, la nueva solución se
evalúa mediante el grafo de restricciones. Esto es necesario
porque un cambio en el orden de una máquina puede introducir
ciclos al combinarse con las precedencias de los trabajos.

La selección aleatoria de movimientos permite explorar
diferentes regiones del espacio de búsqueda.

\subsection{Implementación de la heurística OBA2-N}

Una vez definidos los mecanismos de representación, evaluación
y generación de vecinos, se integró la heurística
\textit{Old Bachelor Acceptance}, en su variante OBA2-N.

El algoritmo mantiene una solución actual, denominada
metafóricamente el ``soltero'', y utiliza un umbral dinámico
para decidir si acepta una solución candidata.

Sea $X$ la solución actual y $X'$ una solución vecina.
El criterio de aceptación por umbral se expresa como:

\[
f(X') < f(X) + T
\]

donde $T$ representa el umbral de aceptación.

Este criterio permite aceptar soluciones que no necesariamente
mejoran el costo de la solución actual, siempre que se
encuentren dentro del límite establecido por el umbral.
Por lo tanto, la búsqueda no se restringe únicamente a
movimientos de mejora.

\subsubsection{Actualización del umbral}

La variante OBA2 utiliza un mecanismo de actualización
dinámica del umbral. Su comportamiento se caracteriza por
un descenso pronunciado y un incremento moderado, lo que
permite modificar el nivel de exigencia durante el proceso
de búsqueda.

El tamaño del decremento aumenta de manera cuadrática,
mientras que el incremento sigue un comportamiento lineal.
Además, se considera la cantidad de aceptaciones fáciles,
relacionadas con el tiempo transcurrido desde la última
aceptación.

La variante OBA2-N incorpora una modificación específica
en el criterio de aceptación: cuando el umbral toma valores
negativos, estos se consideran como cero para aplicar la
regla de aceptación.

Esto permite establecer un límite al efecto de los umbrales
negativos durante la decisión de aceptar una solución
candidata. Es importante distinguir que esta modificación
se aplica al criterio de aceptación y no significa que
el valor almacenado del umbral deba eliminarse o modificarse
automáticamente.

\subsubsection{Proceso de búsqueda}

El procedimiento general de búsqueda de OBA2-N se organiza
en las siguientes etapas:

\begin{enumerate}
    \item Se recibe una solución inicial y se evalúa su costo.
    \item Se genera una solución vecina mediante un movimiento
    sobre el orden de operaciones de una máquina.
    \item Se evalúa la solución candidata, calculando su
    makespan cuando es factible y su penalización cuando
    presenta ciclos.
    \item Se aplica el criterio de aceptación de OBA2-N,
    utilizando el umbral correspondiente.
    \item Si la solución candidata es aceptada, se convierte
    en la nueva solución actual.
    \item Se actualiza el umbral de acuerdo con las reglas
    de la heurística.
    \item Se conserva la mejor solución encontrada durante
    el proceso.
    \item El procedimiento continúa hasta alcanzar el número
    de iteraciones establecido.
\end{enumerate}

Durante la búsqueda resulta conveniente distinguir entre
la solución actual, la mejor solución observada y la mejor
solución factible. Esta separación permite conservar el
mejor calendario válido encontrado, incluso cuando la
solución actual o la mejor solución global no sean
factibles.

\subsection{Inicialización de soluciones}

Antes de ejecutar la heurística es necesario construir un
calendario inicial. Para ello, se consideran las operaciones
de la instancia y se agrupan de acuerdo con las máquinas
que deben procesarlas.

La inicialización aleatoria permite generar diferentes
órdenes de procesamiento para las máquinas. Sin embargo,
mezclar independientemente todas las secuencias puede
producir calendarios con ciclos, debido a que no se
garantiza que las precedencias de los trabajos sean
respetadas.

Por esta razón, la solución inicial debe someterse al
mismo procedimiento de evaluación utilizado durante
la búsqueda, con el propósito de identificar si es
factible y calcular su costo.

La utilización de distintas semillas permite obtener
diferentes soluciones iniciales y explorar distintas
trayectorias de búsqueda.

\subsection{Experimentación y registro de resultados}

Para analizar el comportamiento de la heurística se
considera la ejecución de múltiples corridas independientes
sobre una misma instancia. Cada corrida utiliza una
semilla diferente para controlar las decisiones aleatorias
del algoritmo.

La utilización de un generador aleatorio local por corrida
permite reproducir los experimentos cuando se mantienen
constantes la instancia, la configuración y la semilla.

Durante cada ejecución se recopilan métricas que permiten
analizar el comportamiento de OBA2-N, entre ellas:

\begin{itemize}
    \item Makespan de la mejor solución factible encontrada.
    \item Costo de la solución.
    \item Número de iteraciones ejecutadas.
    \item Cantidad de movimientos aceptados.
    \item Valor final del umbral.
    \item Tiempo de ejecución.
    \item Semilla utilizada.
    \item Indicador de factibilidad de la solución.
\end{itemize}

Los resultados de cada corrida se almacenan en archivos
de texto independientes. Cada archivo contiene la
representación del calendario y la información relacionada
con la ejecución.

Adicionalmente, se contempla la generación de un resumen
global que permita comparar los resultados obtenidos
con las diferentes semillas. Entre las medidas de interés
se encuentran el mejor, peor y promedio del makespan
de las corridas factibles, así como el porcentaje de
ejecuciones que lograron encontrar una solución factible.

\subsection{Flujo general del programa}

La ejecución completa del proyecto sigue una secuencia
de etapas que integra los módulos descritos anteriormente:

\begin{enumerate}
    \item \textbf{Lectura de la instancia:} se recibe el archivo
    de entrada y se construyen los trabajos y operaciones.

    \item \textbf{Validación:} se comprueba la consistencia
    de los datos y se calculan las cotas inferior y superior.

    \item \textbf{Configuración:} se establecen los parámetros
    necesarios para ejecutar OBA2-N.

    \item \textbf{Inicialización:} se genera un calendario
    inicial y se evalúa su factibilidad.

    \item \textbf{Generación de vecinos:} se selecciona una
    máquina y se realiza un movimiento sobre su secuencia
    de operaciones.

    \item \textbf{Evaluación y aceptación:} se calcula el
    costo de la solución candidata y se aplica el criterio
    de aceptación de OBA2-N.

    \item \textbf{Actualización:} se actualiza el umbral y
    se registran las mejores soluciones encontradas.

    \item \textbf{Finalización:} al alcanzar el límite de
    iteraciones, se guardan las métricas y el resultado
    de la corrida.

    \item \textbf{Experimentación:} se repite el procedimiento
    utilizando diferentes semillas para comparar los
    resultados.
\end{enumerate}

La separación de estas etapas permite mantener independiente
la lógica del problema respecto de la heurística de
optimización, facilitando la revisión, prueba y modificación
de los componentes del proyecto.

\subsection{Estructura del proyecto}
\label{sec:estructura}

El proyecto se organizó mediante una estructura modular utilizando
. Esta organización permite separar las responsabilidades relacionadas
con la definición del problema, la representación de soluciones,
la evaluación de calendarios y la ejecución de la heurística OBA2-N.

Como se muestra en la Figura~\ref{fig:arbol}, el proyecto se divide
principalmente en los directorios \texttt{cmd}, \texttt{data},
\texttt{gantt} e \texttt{internal}. Cada uno contiene los archivos
necesarios para realizar una tarea específica dentro del proceso
de resolución del JSSP.

\begin{itemize}

    \item \textbf{cmd:} contiene el punto de entrada de la aplicación.
    Dentro de \texttt{jssp-oba} se encuentra el archivo
    \texttt{main.go}, encargado de coordinar la ejecución del
    programa, recibir los parámetros necesarios y conectar los
    diferentes módulos del proyecto.

    \item \textbf{data:} almacena los archivos utilizados y
    generados durante los experimentos. Se divide en cuatro
    directorios:
    \begin{itemize}
        \item \texttt{configs}: contiene los archivos JSON con
        los parámetros de configuración para las diferentes
        instancias y ejecuciones del algoritmo.

        \item \texttt{instances}: almacena las instancias del
        problema JSSP que serán utilizadas como entrada del
        programa.

        \item \texttt{seeds}: contiene el archivo con las semillas
        utilizadas para controlar las decisiones aleatorias y
        permitir la repetición de los experimentos.

        \item \texttt{results}: almacena los archivos de salida
        generados durante las ejecuciones, incluyendo la
        información de las soluciones y las métricas obtenidas.
    \end{itemize}

    \item \textbf{gantt:} contiene los recursos relacionados
    con la generación de diagramas de Gantt. El archivo
    \texttt{gantt.py} permite procesar la información de una
    solución y generar una representación gráfica del
    calendario de operaciones.

    \item \textbf{internal:} contiene la implementación de los
    módulos principales del proyecto. En este directorio se
    concentra la lógica necesaria para modelar el problema,
    generar soluciones, evaluarlas y ejecutar la heurística
    de optimización. Se divide en los siguientes paquetes:

    \begin{itemize}

        \item \texttt{oba}: implementa la heurística
        \textit{Old Bachelor Acceptance}, variante OBA2-N.
        El archivo \texttt{config.go} contiene los parámetros
        necesarios para configurar el algoritmo, mientras que
        \texttt{oba.go} contiene la lógica del proceso de
        búsqueda y el criterio de aceptación de soluciones.

        \item \texttt{neighborhood}: se encarga de generar
        soluciones vecinas a partir de un calendario existente.
        El archivo \texttt{generate.go} gestiona la generación
        de candidatos y \texttt{moves.go} contiene los
        movimientos que modifican el orden de las operaciones
        dentro de las máquinas.

        \item \texttt{objective}: contiene los procedimientos
        utilizados para evaluar las soluciones. 
        \texttt{cycles.go} construye el grafo de precedencias
        y permite identificar ciclos; \texttt{makespan.go}
        calcula el tiempo total de finalización de los trabajos,
        y \texttt{evaluate.go} integra estas operaciones para
        obtener el valor de la función objetivo.

        \item \texttt{problem}: define las estructuras de datos
        utilizadas para representar una instancia del JSSP.
        \texttt{instance.go} y \texttt{operation.go} contienen
        las estructuras correspondientes a los trabajos y
        operaciones. Por su parte, \texttt{parser.go} permite
        leer las instancias desde archivos, \texttt{validate.go}
        verifica la consistencia de los datos y
        \texttt{bounds.go} calcula las cotas inferior y superior
        del problema.

        \item \texttt{reporting}: contiene los procedimientos
        para almacenar los resultados de las ejecuciones.
        El archivo \texttt{export.go} se encarga de generar
        los archivos de salida con la información de las
        soluciones y sus métricas.

        \item \texttt{solution}: define la representación
        de los calendarios utilizados durante la búsqueda.
        \texttt{schedule.go} contiene la estructura principal
        de una solución, \texttt{representation.go} proporciona
        las operaciones relacionadas con su representación
        y \texttt{clone.go} permite crear copias independientes
        de los calendarios antes de aplicar modificaciones.

        \item \texttt{solver}: coordina la preparación y
        ejecución de los experimentos. 
        \texttt{initialize.go} construye la solución inicial,
        \texttt{seeds.go} gestiona las semillas utilizadas
        para las ejecuciones aleatorias y
        \texttt{experiment.go} permite organizar las corridas
        experimentales.

    \end{itemize}

    \item \textbf{go.mod:} define el módulo de Go y administra
    la información necesaria para compilar y organizar las
    dependencias del proyecto.

    \item \textbf{README.md:} contiene la documentación general
    del proyecto, incluyendo información sobre su propósito,
    configuración y ejecución.

\end{itemize}

Esta organización facilita el mantenimiento del código,
la realización de pruebas individuales y la posibilidad de
modificar algún componente del algoritmo sin afectar
directamente al resto de la implementación.

\begin{figure}[H]
    \centering
    \begin{minipage}{0.92\textwidth}
\begin{lstlisting}
.
|-- cmd
|   |-- jssp-oba
|       |-- main.go
|-- data
|   |-- configs
|   |   |-- config100x100.json
|   |   |-- config50x20.json
|   |-- instances
|   |   |-- config100x100.json
|   |   |-- config40x40.json
|   |   |-- config50x20.json
|   |-- seeds
|   |   |-- seeds.txt
|   |-- results
|       |-- resultado_20261106_133135.368687566.txt
|-- gantt
|   |-- gantt.png
|   |-- gantt.py
|-- go.mod
|-- internal
|   |-- oba
|   |   |-- config.go
|   |   |-- oba.go
|   |   |-- oba.go
|   |-- neighborhood
|   |   |-- generate.go
|   |   |-- moves.go
|   |-- objective
|   |   |-- cycles.go
|   |   |-- evaluate.go
|   |   |-- makespan.go
|   |-- problem
|   |   |-- bounds.go
|   |   |-- instance.go
|   |   |-- operation.go
|   |   |-- parser.go
|   |   |-- validate.go
|   |-- reporting
|   |   |-- export.go
|   |-- solution
|   |   |-- clone.go
|   |   |-- representation.go
|   |   |-- schedule.go
|   |-- solver
|   |   |-- experiment.go
|   |   |-- initialize.g
|   |   |-- seeds.go
|-- README.md
\end{lstlisting}
    \end{minipage}
    \caption{Estructura de directorios del proyecto.}
    \label{fig:arbol}
\end{figure}


\section{Resultados}


% =====================================================
% 3. CONCLUSIÓN
% =====================================================
\section{Conclusión}


% =====================================================
% 4. REFERENCIAS
% =====================================================
\begin{thebibliography}{99}

\bibitem{OBA}
Huy, T. C., Kahng, A. B. y Tsao, C. W. Old Bachelor Acceptance: A New Class of Non-Monotone Threshold Accepting Methods.

\bibitem{JSSP}
Aguilar, D. J. V. (2026). JSSP - Desarrollo del Problema. Documento de especificación del problema.

\end{thebibliography}

\end{document}
